\heading{实验物理中的统计方法-第11章 解谱法}


\begin{question}
考虑图 9.1 的探测器设备。假设 \(x\) 的分辨率函数服从高斯分布

\[
s(x|x')=\frac{1}{\sqrt{2\pi}\sigma_x}\exp\left[-\frac{(x-x')^2}{2\sigma_x^2}\right].
\]

求 \(\cos\theta=a/\sqrt{x^2+a^2}\) 的分辨率函数。
\begin{solution}
记
\[
c=\cos\theta=\frac{a}{\sqrt{x^2+a^2}},
\qquad 0<c\le1.
\]
由于这个变换依赖于 \(x^2\)，一般不是一一对应的。对给定 \(c\)，反解为
\[
x_\pm(c)=\pm a\frac{\sqrt{1-c^2}}{c}.
\]
并且
\[
\left|\frac{dx_\pm}{dc}\right|
=\frac{a}{c^2\sqrt{1-c^2}}.
\]

若 \(x\) 是可正可负的测量坐标，则测得的 \(c\) 在真值 \(x'\) 下的分辨率函数应对两个分支求和：
\[
\begin{split}
s_c(c\mid x')
&=\sum_{\epsilon=\pm1}
 s\left(\epsilon a\frac{\sqrt{1-c^2}}{c}\middle|x'\right)
 \frac{a}{c^2\sqrt{1-c^2}}\\
&=\frac{a}{c^2\sqrt{1-c^2}}
\frac{1}{\sqrt{2\pi}\sigma_x}
\left\{
\exp\left[-\frac{(x_+(c)-x')^2}{2\sigma_x^2}\right]
+
\exp\left[-\frac{(x_-(c)-x')^2}{2\sigma_x^2}\right]
\right\}.
\end{split}
\]
其中 \(0<c\le1\)。

若问题物理上只允许 \(x\ge0\) 的单支（例如 \(x\) 表示非负距离），则只保留正分支：
\[
s_c(c\mid c')=
\frac{1}{\sqrt{2\pi}\sigma_x}
\exp\left[-\frac{(x(c)-x(c'))^2}{2\sigma_x^2}\right]
\frac{a}{c^2\sqrt{1-c^2}},
\]
其中
\[
x(c)=a\frac{\sqrt{1-c^2}}{c}.
\]
\pyfig[0.72\linewidth]{figures/fig_11_01_cos_resolution.png}{单支 \(x\ge0\) 情形下，非线性变换后 \(c=\cos\theta\) 分辨率函数的示例。}
\end{solution}
\end{question}


\begin{question}
对等区间宽度的直方图 \(\mu=(\mu_1,\ldots,\mu_M)\)，考虑 \(k=1\) 的
Tikhonov 正规化函数，

\[
S(\mu)=-\sum_{i=1}^{M-1}(\mu_i-\mu_{i+1})^2.
\]

求 \(M\times M\) 维矩阵 \(G\)，使得 \(S(\mu)\) 可以表示成以下形式

\[
S(\mu)=-\sum_{i,j=1}^{M}G_{ij}\mu_i\mu_j=-\mu^T G\mu.
\]
\begin{solution}
展开平方和：
\[
\sum_{i=1}^{M-1}(\mu_i-\mu_{i+1})^2
=\mu_1^2+\mu_M^2+2\sum_{i=2}^{M-1}\mu_i^2
-2\sum_{i=1}^{M-1}\mu_i\mu_{i+1}.
\]
因此 \(S(\mu)=-\mu^T G\mu\)，其中
\[
G=
\begin{pmatrix}
1 & -1 & 0 & \cdots & 0\\
-1 & 2 & -1 & \ddots & \vdots\\
0 & -1 & 2 & \ddots & 0\\
\vdots & \ddots & \ddots & 2 & -1\\
0 & \cdots & 0 & -1 & 1
\end{pmatrix}.
\]
等价地，
\[
G_{11}=G_{MM}=1,\qquad
G_{ii}=2\ (2\le i\le M-1),\qquad
G_{i,i+1}=G_{i+1,i}=-1,
\]
其余矩阵元为 \(0\)。这就是离散一阶差分平方对应的二次型矩阵。
\pyfig[0.56\linewidth]{figures/fig_11_02_tikhonov_matrix.png}{\(M=8\) 时一阶 Tikhonov 矩阵 \(G\) 的结构示意。}
\end{solution}
\end{question}


\begin{question}
考虑期待值为 \(\mu=(\mu_1,\ldots,\mu_M)\) 的直方图，对应的概率为
\(p=\mu/\mu_{\rm tot}\)，其中 \(\mu_{\rm tot}=\sum_{i=1}^{M}\mu_i\)。

\begin{enumerate}
\item
  证明 Shannon 熵
\end{enumerate}

\[
H(p)=-\sum_{i=1}^{M}p_i\log p_i,
\]

在所有区间 \(i\) 的 \(p_i=1/M\) 时最大。（利用 Lagrange
乘子法引入限制条件 \(\sum_{i=1}^{M}p_i=1\)。）

\begin{enumerate}[start=2]
\item
  证明交叉熵
\end{enumerate}

\[
K(p;q)=-\sum_{i=1}^{M}p_i\log\frac{p_i}{Mq_i},
\]

在概率 \(p\) 等于参考分布 \(q\) 时最大。
\begin{solution}
设 \(p_i=\mu_i/\mu_{\rm tot}\)，于是 \(\sum_i p_i=1\)。

\begin{enumerate}
\item 在约束 \(\sum_i p_i=1\) 下，对 Shannon 熵
\[
H(\mathbf p)=-\sum_{i=1}^M p_i\log p_i
\]
作 Lagrange 乘子法：
\[
L(\mathbf p,\lambda)
=-\sum_i p_i\log p_i+\lambda\left(\sum_i p_i-1\right).
\]
由
\[
\frac{\partial L}{\partial p_i}=-(\log p_i+1)+\lambda=0
\]
可知所有 \(p_i\) 相等。结合归一化条件，得到
\[
p_i=\frac1M,\qquad H_{\max}=\log M.
\]

\item 将题中的量改写为
\[
K(\mathbf p;\mathbf q)
=-\sum_{i=1}^M p_i\log\frac{p_i}{Mq_i}
=\log M-\sum_{i=1}^M p_i\log\frac{p_i}{q_i}
=\log M-D_{\mathrm{KL}}(\mathbf p\Vert\mathbf q).
\]
由于 \(D_{\mathrm{KL}}(\mathbf p\Vert\mathbf q)\ge 0\)，且等号仅在
\(\mathbf p=\mathbf q\) 时成立，所以
\[
K(\mathbf p;\mathbf q)\le \log M.
\]
因此 \(K\) 在 \(p=q\) 时达到最大值。
\end{enumerate}
\end{solution}
\end{question}


\begin{question}
考虑观测到的直方图 \(n=(n_1,\ldots,n_N)\)，对应的期待值
\(\nu=(\nu_1,\ldots,\nu_N)\) 与真值直方图 \(\mu=(\mu_1,\ldots,\mu_N)\)
的关系为 \(\nu=R\mu\)。假设协方差矩阵 \(V\) 和响应矩阵 \(R\)
已知，并且直方图中没有本底。

\begin{enumerate}
\item
  通过最大化 \(\Phi(\mu)\) 构造 \(\mu\) 的估计量
\end{enumerate}

\[
\Phi(\mu)=-\frac{\alpha}{2}\chi^2(\mu)+S(\mu)
\]

\[
=-\frac{\alpha}{2}(n-R\mu)^T V^{-1}(n-R\mu)-\mu^T G\mu,
\]

其中 \(\alpha\) 为正规化参数，\(M\times M\) 对称矩阵 \(G\)
由已知常数确定（参考 Statistical Data Analysis 中 11.5.1
节）。证明估计量 \(\hat\mu\) 为

\[
\hat\mu=(\alpha R^T V^{-1}R+2G)^{-1}\alpha R^T V^{-1}n,
\]

并求协方差矩阵 \(U_{ij}=\operatorname{cov}[\hat\mu_i,\hat\mu_j]\)。

\begin{enumerate}[start=2]
\item
  考虑限制条件
  \(\nu_{\rm tot}=\sum_{i=1}^{N}\nu_i=\sum_{i=1}^{N}\sum_{j=1}^{M}R_{ij}\mu_j\)
  等于总观测事例数 \(n_{\rm tot}=\sum_{i=1}^{N}n_i\)，通过对参数 \(\mu\)
  和 Lagrange 乘子 \(\lambda\) 最大化 \(\phi(\mu)\) 求得结果
\end{enumerate}

\[
\phi(\mu)=-\frac{\alpha}{2}(n-R\mu)^T V^{-1}(n-R\mu)-\mu^T G\mu+\lambda(n_{\rm tot}-\nu_{\rm tot}).
\]

求估计量 \(\hat\mu\) 及其协方差。

\begin{enumerate}[start=3]
\item
  利用 Statistical Data Analysis 中方程 (11.76)，对 (a) 和 (b)
  两种情况分别构造偏置 \(b=E[\hat\mu]-\mu\) 的估计量 \(\hat b\)。
\end{enumerate}
\begin{solution}
记
\[
M=\alpha R^T V^{-1}R+2G.
\]

\begin{enumerate}
\item 无总数约束

对
\[
\Phi(\mu)=
-\frac\alpha2(n-R\mu)^TV^{-1}(n-R\mu)-\mu^TG\mu
\]
求驻点：
\[
\alpha R^TV^{-1}(n-R\hat\mu)-2G\hat\mu=0.
\]
故
\[
\hat\mu=M^{-1}\alpha R^TV^{-1}n.
\]
令
\[
A=M^{-1}\alpha R^TV^{-1},
\]
则 \(\hat\mu=An\)。若 \(\operatorname{cov}(n)=V\)，则
\[
U=\operatorname{cov}(\hat\mu)=AVA^T
=\alpha^2M^{-1}R^TV^{-1}R M^{-1}.
\]

\item 加总数约束

令 \(\mathbf 1_N\) 为 \(N\) 维全 1 向量，
\[
c=R^T\mathbf 1_N.
\]
约束为
\[
c^T\mu=\mathbf 1_N^Tn.
\]
带 Lagrange 乘子的驻点方程为
\[
M\hat\mu=\alpha R^TV^{-1}n+\lambda c.
\]
先定义无约束线性解
\[
\mu_0=M^{-1}\alpha R^TV^{-1}n.
\]
由约束得
\[
\lambda=
\frac{\mathbf 1_N^Tn-c^T\mu_0}{c^TM^{-1}c}.
\]
所以
\[
\hat\mu
=\mu_0+M^{-1}c\frac{\mathbf 1_N^Tn-c^T\mu_0}{c^TM^{-1}c}.
\]
为避免排版歧义，上式也可写成
\[
\hat\mu=Bn,
\]
其中
\[
B=M^{-1}\alpha R^TV^{-1}
+M^{-1}c\,
\frac{\mathbf 1_N^T-c^TM^{-1}\alpha R^TV^{-1}}{c^TM^{-1}c}.
\]
于是
\[
U_{\rm constr}=BVB^T.
\]

\item 偏置估计
\pyfig[0.58\linewidth]{figures/fig_11_04_response_matrix.png}{解谱问题中一个响应矩阵的示意；列归一化表示每个真值 bin 对观测 bin 的迁移概率。}
若线性估计量写成 \(\hat\mu=An\)，且 \(E[n]=R\mu\)，则
\[
b=E[\hat\mu]-\mu=(AR-I)\mu.
\]
实际中 \(\mu\) 未知，按书中式 (11.76) 的思想可用估计量代入，取
\[
\hat b=(AR-I)\hat\mu.
\]
无约束情形
\[
A=M^{-1}\alpha R^TV^{-1},
\]
并且利用 \(M=\alpha R^TV^{-1}R+2G\)，有
\[
AR-I=-2M^{-1}G,
\]
因此
\[
\hat b_a=-2M^{-1}G\hat\mu.
\]

带约束情形只需把 \(A\) 换成 \(B\)：
\[
\hat b_b=(BR-I)\hat\mu.
\]
若展开，可得
\[
BR-I=-2M^{-1}
\left(
G-\frac{c\,c^TM^{-1}G}{c^TM^{-1}c}
\right),
\]
所以
\[
\hat b_b=-2M^{-1}
\left(
G-\frac{c\,c^TM^{-1}G}{c^TM^{-1}c}
\right)\hat\mu.
\]
\end{enumerate}
\end{solution}
\end{question}

